\capitulo{4}{Técnicas y herramientas}
En este apartado se va a realizar una breve explicación de la metodología empleada, asi como de las herramientas utilizadas para el desarrollo del proyecto. En dicha explicación se va a incluir una comparación con otras herramientas similares en la cual se muestre el motivo de la selección de dicha herramienta. 
\section{Metodologías}

\subsection{Scrum}
La metodología Scrum es un marco de trabajo ágil orientado a la gestión y desarrollo de proyectos de manera iterativa e incremental. Se organiza en ciclos cortos denominados \textit{sprints}, al final de los cuales se entrega un incremento funcional del producto. Estos \textit{sprints} favorecen la mejora continua y la adaptación al cambio, además, aseguran una entrega temprana de valor. A diferencia de metodologías tradicionales, Scrum aporta mayor flexibilidad y retroalimentación constante~\cite{scrum:guide}. 

\section{Repositorio}
Para alojar el repositorio se han valorado distintas alternativas como \href{https://bitbucket.org/}{Bitbucket}, \href{https://github.com/}{GitHub}, \href{https://gitlab.com/}{GitLab} o \href{https://trello.com/}{Trello}. Finalmente se ha elegido GitHub que es una plataforma basada en la nube que permite alojar repositorios de código. GitHub ofrece funciones muy útiles como ramas (\textit{branches}), \textit{pull request}, seguimiento de cambios (\textit{commits}), historial del proyecto, gestión de incidencias (\textit{issues}). Además, permite la colaboración entre usuarios, lo que es muy útil para la revisión del proyecto.

Otra ventaja es que GitHub se integra con herramientas de integración y despliegue continuo, lo que facilita la automatización de pruebas y la entrega de \textit{software}. Asimismo, permite incluir documentación directamente en el repositorio, ya sea mediante archivos \textit{README} o a través de \textit{Wikis}, lo que mejora la claridad y la reproducibilidad del trabajo. Finalmente, al tratarse de la plataforma más utilizada a nivel mundial, cuenta con abundante documentación, lo que la convierte en una alternativa sólida y con amplio soporte~\cite{gitHub}.

\section{Gestión de tareas}
Para la gestión de las tareas del proyecto se han valorado diversas herramientas, como \href{https://www.zenhub.com/}{ZenHub}, \href{https://github.com/}{GitHub
  Projects} o \href{https://asana.com/es}{Asana}. Al final se ha optado por usar ZenHub que es una herramienta de gestión de proyectos fácilmente integrable con GitHub. Se ha elegido porque tiene mayor funcionalidad que GitHub Proyects. Ambos permiten hacer tableros Kanban para organizar las tareas (\textit{to-do, in progress, doing}). Pero, ZenHub, además, permite asignar \textit{story points} a las tareas y generar los gráficos \textit{burndown} de los \textit{sprints}.

\section{\LaTeX{}}
Para la edición de los documentos en \LaTeX{} se ha considerado usar \href{http://www.xm1math.net/texmaker/}{\TeX{Maker}}, \href{https://www.texstudio.org/}{\TeX{Studio}}, \href{https://code.visualstudio.com/}{Visual Studio Code} y \href{https://es.overleaf.com/}{Overleaf}. Finalmente se ha optado por usar \TeX{Maker}, ya que es un editor multiplataforma, gratuito y de código abierto. Tiene un visor PDF integrado, así como herramientas para gestionar la bibliografía mediante Bib\TeX{}. Además, tiene una interfaz simple que facilita el aprendizaje, incorpora funciones como el autocompletado, el plegado de código y la detección de errores de compilación, sin necesidad de instalar complementos adicionales.

Como distribución \LaTeX{} para procesar los documentos .tex se ha valorado el uso de \href{https://miktex.org/}{Mik\TeX{}} y \href{https://www.tug.org/texlive/}{\TeX{Live}}, al final se ha elegido usar Mik\TeX{} ya que permite realizar una instalación mínima e ir añadiendo automáticamente los paquetes necesarios durante la compilación. Además, presenta un buen soporte para Windows, lo que se adapta bien a mi entorno de trabajo~\cite{miktex:textlive}.

\section{\textit{Web Scraping}}
Para el \textit{web scraping} se ha valorado usar \href{https://selenium-python.readthedocs.io/}{Selenium} y \href{https://playwright.dev/}{Playwright}.
Selenium controla navegadores reales y es un estándar veterano, pero Playwright es más completo. Ya que ofrece automatización multi-navegador con contextos aislados, funcionamiento headless, espera automática de elementos y APIs consistentes, lo que lo vuelve más robusto y rápido para páginas con JavaScript pesado y flujos de login. En la tabla~\ref{tabla:SeleniumPlaywright} se muestran las ventajas y desventajas de cada uno.

\begin{table}[!htbp]
	\scriptsize 
	\centering
	\renewcommand{\arraystretch}{1.1}
	\begin{tabularx}{\linewidth}{
    		>{\centering\arraybackslash}p{0.18\linewidth}
    		>{\centering\arraybackslash}X
    		>{\centering\arraybackslash}X}
		\toprule
			& \textbf{Ventajas} & \textbf{Desventajas} \\
		\midrule
		\textbf{\textit{Selenium}}~\cite{selenium:docs} & 
		
			\begin{itemize}
			\tightlist
				\item Utiliza el estándar W3C \textit{WebDriver}~\cite{w3c:webdriver2}, lo que aporta interoperabilidad y soporte a largo plazo.
				\item Multilenguaje y ecosistema maduro. Permite una integración fácil.
				\item Es robusto para escalado distribuido. 
			\end{itemize}  
			& 			
			 \begin{itemize}
			\tightlist
				\item Es propenso a fallos (\textit{flakiness}) si no se dominan las esperas.
				\item La intercepción de red en menos homogénea entre navegadores.
			\end{itemize}
			
			\\
			
		\textbf{\textit{Playwright}}~\cite{playwright:actionability} & 
		 
		 \begin{itemize}
			\tightlist
				\item Es muy estable y rápido en \textit{webs} con \textit{JavaScript} pesado.
				\item Presenta multitud de herramientas integradas como \textit{Trace Viewer}, capturas de pantalla y videos.
				\item Contiene selectores potentes y permite manejar los \textit{inframes} de forma sencilla. 
				\item Tiene un control de contexto aislado muy cómodo.
				\item Su modelos asíncrono permite realizar acciones concurrentes sin bloquear el hilo empleando la \textit{API async/await}.
			\end{itemize}   
			& 			
			\begin{itemize}
			\tightlist
				\item No utiliza el estándar \textit{WebDriver}. 
			\end{itemize}   
			
			\\
		
		\bottomrule
	\end{tabularx}
	\caption{Ventajas y desventajas de \textit{Selenium} y \textit{Playwright}.}
	\label{tabla:SeleniumPlaywright}
\end{table}

Aunque \textit{Selenium} sigue siendo una apuesta segura, para \textit{scraping} moderno con multipes sesiones y depuración rápida, \textit{Playwright} requiere menos código y es más robusto. 

\section{Gestión de dependencias y entornos \textit{Python}}
El proyecto requiere entornos aislados y reproducibles que faciliten el trabajo local, la integración continua y la generación de entregables. Se ha valorado utilizar \href{http://python-poetry.org/docs/}{\textit{Poetry}} y \href{https://docs.astral.sh/uv/}{\textit{uv}}. \textit{Poetry} es una herramienta madura de gestión de dependencias y empaquetado que centraliza la configuración en \textit{pyproject.toml}. Crea y gestiona entornos virtuales de forma automática, mantiene un archivo de bloqueo (\textit{poetry.lock}) para asegurar instalaciones deterministas, permite definir grupos de dependencias, así como, generar de manera automática el \textit{requirements.txt}. Por su parte \textit{uv} es un gestor rápido que, al igual que \textit{Poetry}, tiene soporte para \textit{pyproject.toml}. Posee un buen rendimiento y una elevada velocidad gracias al uso de \textit{caches} eficientes. Finalmente, se ha optado por usar \textit{Poetry} por su cobertura integral del ciclo de vida dentro de un flujo coherente y estable.